\section{USAMO 2017}
        \begin{enumerate}
        	\item (\textit{USAMO 2017}) Note: For any geometry problem whose statement begins with an asterisk ($*$), the first page of the solution must be a large, in-scale, clearly labeled diagram. Failure to meet this requirement will result in an automatic 1-point deduction. 


 Prove that there are infinitely many distinct pairs $(a,b)$ of relatively prime positive integers $a > 1$ and $b > 1$ such that $a^b + b^a$ is divisible by $a + b.$


 



\textbf{Solution 1}

Let $n=a+b$. Since $gcd(a,b)=1$, we know $gcd(a,n)=1$. We can rewrite the condition as 

 
 \[a^{n-a}+(n-a)^a \equiv 0 \mod{n}\]
 \[a^{n-a}\equiv-(-a)^a \mod{n}\]Assume $a$ is odd. Since we need to prove an infinite number of pairs exist, it suffices to show that infinitely many pairs with $a$ odd exist.

 Then we have 
 \[a^{n-a}\equiv a^a \mod{n}\]
 \[1 \equiv a^{2a-n} \mod{n}\]

 We know by Euler's theorem that $a^{\varphi(n)} \equiv 1 \mod{n}$, so if $2a-n=\varphi(n)$ we will have the required condition.

 This means $a=\frac{n+\varphi(n)}{2}$. Let $n=2p$ where $p$ is a prime, $p\equiv 1\mod{4}$. Then $\varphi(n) = 2p*\left(1-\frac{1}{2}\right)\left(1-\frac{1}{p}\right) = p-1$, so 
 \[a = \frac{2p+p-1}{2} = \frac{3p-1}{2}\]Note the condition that $p\equiv 1\mod{4}$ guarantees that $a$ is odd, since $3p-1 \equiv 2\mod{4}$

 This makes $b = \frac{p+1}{2}$. Now we need to show that $a$ and $b$ are relatively prime. We see that
 \[gcd\left(\frac{3p-1}{2},\frac{p+1}2\right)=\frac{gcd(3p-1,p+1)}{2}\]
 \[=\frac{gcd(p-3,4)}{2}=\frac22=1\]By the Euclidean Algorithm.

 Therefore, for all primes $p \equiv 1\mod{4}$, the pair $\left(\frac{3p-1}{2},\frac{p+1}{2}\right)$ satisfies the criteria, so infinitely many such pairs exist.

 



\textbf{Solution 2}

Take $a=2n-1, b=2n+1, n\geq 2$. It is obvious (use the Euclidean Algorithm, if you like), that $\gcd(a,b)=1$, and that $a,b>1$.

 Note that

 
 \[a^2 = 4n^2-4n+1 \equiv 1 \pmod{4n}\]

 
 \[b^2 = 4n^2+4n+1 \equiv 1 \pmod{4n}\]

 So

 
 \[a^b+b^a = a(a^2)^n+b(b^2)^{n-1} \equiv\] 
 \[a\cdot 1^n + b\cdot 1^{n-1} \equiv a+b = 4n \equiv 0 \pmod{4n}\]

 Since $a+b=4n$, all such pairs work, and we are done.

 



\textbf{Solution 3}

Let $x$ be odd where $x>1$. We have $x^2-1=(x-1)(x+1),$ so $x^2-1 \equiv 0 \pmod{2x+2}.$ This means that  $x^{x+2}-x^x \equiv 0 \pmod{2x+2},$ and since x is odd, $x^{x+2}+(-x)^x \equiv 0 \pmod{2x+2},$ or $x^{x+2}+(x+2)^x \equiv 0 \pmod{2x+2},$ as desired.

 



\textbf{Solution 4}

I claim that the ordered pair $(2^{n} - 1, 2^{n} + 1)$ satisfies the criteria for all $n \geq 2.$ Proof:It is easy to see that the order modulo $2^{n+1}$ of $(2^{n} - 1)^k$ is $2,$ since $(2^{n} - 1)^2 = 2^{2n} - 2 \cdot 2^{n} + 1 \equiv 1 \mod 2^{n+1}.$, and we can assert and prove similarly for the order modulo $2^{n+1}$ of $(2^{n} + 1)^k.$ Thus, the remainders modulo $2^{n+1}$ that the sequences of powers of $2^{n} - 1$ and $2^{n} + 1$ generate are $1$ for even powers and $2^{n} - 1$ and $2^{n} + 1$ for odd powers, respectively. Since $2^{n} + 1$ and $2^{n} - 1$ are both odd for $n \geq 2,$ 
 \[(2^n + 1)^{2^n - 1} + (2^n - 1)^{2^n + 1} \equiv 0 \mod 2^{n+1}.\] Since there are infinitely many powers of $2$ and since all ordered pairs $(2^n - 1, 2^n+1)$ contain relatively prime integers, we are done.$\boxed{}$

 -fidgetboss\_4000

 



\textbf{Solution 5 (Motivation for Solution)}

Note that 
 \[a^b+b^a=a^b-a^a+a^a+b^a.\] To get rid of the $a^a+b^a$ part $\pmod{a+b},$ we can use the sum of powers factorization. However, $a$ must be odd for us to do this. If we assume that $a$ is odd,
 \[a^b-a^a+a^a+b^a\equiv a^b-a^a+(a+b)(\text{an integer})\equiv a^b-a^a \equiv a^a\left(a^{b-a}-1\right)\pmod{a+b}.\]Because $a$ and $b$ are relatively prime, $a+b$ cannot divide $a^a.$ Thus, we have to show that there exists an integer $b$ such that for odd $a,$
 \[a^{b-a}\equiv 1 \pmod{a+b}.\]Suppose that $a=2n-1.$ To keep the powers small, we try $b=2n+k$ for small values of $k.$ We can find that $b=2n$ does not work. $b=2n+1$ works though, as $a+b=4n$ and
 \[(2n-1)^{2n+1-2n+1}\equiv (2n-1)^2\equiv 4n^2-4n+1\equiv 4n(n-1)+1 \equiv 1 \pmod{4n}.\]Because $a$ is odd, $b=a+2$ is relatively prime to $a.$ Thus, 
 \[(a,b)=(2n-1,2n+1)\] is a solution for positive $n\ge 2.$ There are infinitely many possible values for $n,$ so the proof is complete. $\blacksquare$ 

 


	\item (\textit{USAMO 2017}) Note: For any geometry problem whose statement begins with an asterisk ($*$), the first page of the solution must be a large, in-scale, clearly labeled diagram. Failure to meet this requirement will result in an automatic 1-point deduction. 


 Let $m_1, m_2, \ldots, m_n$ be a collection of $n$ positive integers, not necessarily distinct. For any sequence of integers $A = (a_1, \ldots, a_n)$ and any permutation $w = w_1, \ldots, w_n$ of $m_1, \ldots, m_n$, define an $A$-inversion of $w$ to be a pair of entries $w_i, w_j$ with $i < j$ for which one of the following conditions holds:

 \[a_i \ge w_i > w_j,\]

 \[w_j > a_i \ge w_i,\] or

 \[w_i > w_j > a_i.\]
Show that, for any two sequences of integers $A = (a_1, \ldots, a_n)$ and $B = (b_1, \ldots, b_n)$, and for any positive integer $k$, the number of permutations of $m_1, \ldots, m_n$ having exactly $k$ $A$-inversions is equal to the number of permutations of $m_1, \ldots, m_n$ having exactly $k$ $B$-inversions.


 


	\item (\textit{USAMO 2017}) Note: For any geometry problem whose statement begins with an asterisk ($*$), the first page of the solution must be a large, in-scale, clearly labeled diagram. Failure to meet this requirement will result in an automatic 1-point deduction. 


 ($*$) Let $ABC$ be a scalene triangle with circumcircle $\Omega$ and incenter $I$. Ray $AI$ meets $\overline{BC}$ at $D$ and meets $\Omega$ again at $M$; the circle with diameter $\overline{DM}$ cuts $\Omega$ again at $K$. Lines $MK$ and $BC$ meet at $S$, and $N$ is the midpoint of $\overline{IS}$. The circumcircles of  $\triangle KID$ and $\triangle MAN$ intersect at points $L_1$ and $L_2$. Prove that $\Omega$ passes through the midpoint of either $\overline{IL_1}$ or $\overline{IL_2}$.


 



\textbf{Solution}


                    \begin{figure}[H]
                        \centering
                        \includegraphics[width=0.5\textwidth]{Solution/USAMO_I_2017_3_Solution_Figure_1.png}
                    \end{figure}
                    
                    \begin{figure}[H]
                        \centering
                        \includegraphics[width=0.5\textwidth]{Solution/USAMO_I_2017_3_Solution_Figure_2.png}
                    \end{figure}
                    
                    \begin{figure}[H]
                        \centering
                        \includegraphics[width=0.5\textwidth]{Solution/USAMO_I_2017_3_Solution_Figure_3.png}
                    \end{figure}
                    Let $X$ be the point on circle $\Omega$ opposite $M$. This means $\angle MAX = 90^\circ, BC \perp XM.$

 $\angle XKM = \angle DKM = 90^\circ \implies$ the points $X, D,$ and $K$ are collinear.

 Let $D' = BC \cap XM \implies DD' \perp XM \implies$

 $S$ is the orthocenter of $\triangle DMX \implies$ the points $X, A,$ and $S$ are collinear.

 Let $\omega$ be the circle centered at $S$ with radius $R = \sqrt {SK \cdot SM}.$

 We denote $I_\omega$ inversion with respect to $\omega.$

 Note that the circle $\Omega$ has diameter $MX$ and contain points $A, B, C,$ and $K.$

 $I_\omega (K) = M \implies$ circle $\Omega \perp \omega \implies C = I_\omega (B), X = I_\omega (A).$ 

 $I_\omega (K) = M \implies$ circle $KMD \perp \omega \implies D' = I_\omega (D) \in KMD \implies$$\angle DD'M = 90^\circ \implies$  the points $X, D',$ and $M$ are collinear.

 Let $F \in AM, MF = MI.$ It is well known that $MB = MI = MC \implies$

 $\Theta = BICF$ is circle  centered at $M.$ $C = I_\omega (B) \implies \Theta \perp \omega.$

 Let $I' =  I_\omega (I ) \implies I' \in \Theta \implies \angle II'M =  90^\circ.$ $I' =  I_\omega (I ), X =  I_\omega (A ) \implies AII'X$ is cyclic.

 $\angle XI'I = \angle XAI =  90^\circ \implies$  the points $X, I' ,$ and $F$ are collinear.

 $I'IDD'$ is cyclic $\implies \angle I'D'M = \angle I'D'C + 90^\circ =  \angle I'ID + 90^\circ,$$\angle XFM = \angle I'FI = 90^\circ – \angle I'IF = 90^\circ – \angle I'ID  \implies$

 $\angle XFM +  \angle I'D'M = 180^\circ \implies I'D'MF$ is cyclic.

 Therefore point $F$ lies on  $I_\omega (IDK).$

 $FA \perp SX, SI' \perp FX \implies I$ is orthocenter of $\triangle FSX.$

 $N$ is midpoint $SI, M$ is midpoint $FI, I$ is orthocenter of $\triangle FSX, A$ is root of height $FA \implies AMN$ is the nine-point circle of $\triangle FSX \implies I' \in AMN.$ 

 Let $N' = I_\omega (N) \implies R^2 = SN \cdot SN' = SI \cdot  SI' \implies$
 \[\frac {SN'}{SI'} = \frac {SI}{SN} =2 \implies\]$\angle XN'I' = \angle XSI' = 90^\circ – \angle AXI' = \angle IFX \implies N'XIF$ is cyclic. 

 Therefore point $F$ lies on  $I_\omega (AMN) \implies I_\omega(F) = L \implies$

 The points $F, L,$ and $S$ are collinear, $AXFL$ is cyclic.

 Point $I$ is orthocenter $\triangle FSX \implies XI \perp SF, \angle ILS = \angle SI'F = 90^\circ$$\implies$ The points $X, I, E,$ and $L$ are collinear.

 $AXFL$ is circle $\implies AI \cdot IF = IL \cdot XI\implies$

 $AI \cdot \frac {IF}{2} = \frac {IL}{2} \cdot IX \implies AI \cdot IM = EI \cdot IX \implies AEMX$ is cyclic.
 \[E \in \Omega.\]\textbf{vladimir.shelomovskii@gmail.com, vvsss}

 \textbf{Contact} Contact v\_Enhance at \href{https://www.facebook.com/v.Enhance}{https://www.facebook.com/v.Enhance}.

 


	\item (\textit{USAMO 2017}) Note: For any geometry problem whose statement begins with an asterisk ($*$), the first page of the solution must be a large, in-scale, clearly labeled diagram. Failure to meet this requirement will result in an automatic 1-point deduction. 


 Let $P_1$, $P_2$, $\dots$, $P_{2n}$ be $2n$ distinct points on the unit circle $x^2+y^2=1$, other than $(1,0)$. Each point is colored either red or blue, with exactly $n$ red points and $n$ blue points. Let $R_1$, $R_2$, $\dots$, $R_n$ be any ordering of the red points. Let $B_1$ be the nearest blue point to $R_1$ traveling counterclockwise around the circle starting from $R_1$. Then let $B_2$ be the nearest of the remaining blue points to $R_2$ travelling counterclockwise around the circle from $R_2$, and so on, until we have labeled all of the blue points $B_1, \dots, B_n$. Show that the number of counterclockwise arcs of the form $R_i \to B_i$ that contain the point $(1,0)$ is independent of the way we chose the ordering $R_1, \dots, R_n$ of the red points.


 



\textbf{Solution}

I define a sequence to be, starting at $(1,0)$ and tracing the circle counterclockwise, and writing the color of the points in that order - either R or B. For example, possible sequences include $RB$, $RBBR$, $BBRRRB$, $BRBRRBBR$, etc.Note that choosing an $R_1$ is equivalent to choosing an $R$ in a sequence, and $B_1$ is defined as the $B$ closest to $R_1$ when moving rightwards. If no $B$s exist to the right of $R_1$, start from the far left. For example, if I have the above example $RBBR$, and I define the 2nd $R$ to be $R_1$, then the first $B$ will be $B_1$. Because no $R$ or $B$ can be named twice, I can simply remove $R_1$ and $B_1$ from my sequence when I choose them. I define this to be a move. Hence, a possible move sequence of $BBRRRB$ is:$BBR_1RRB_1\implies B_2BRR_2\implies B_3R_3$

 Note that, if, in a move, $B_n$ appears to the left of $R_n$, then $\stackrel{\frown}{R_nB_n}$ intersects $(1,0)$

 Now, I define a commencing $B$ to be a $B$ which appears to the left of all $R$s, and a terminating $R$ to be a $R$ which appears to the right of all $B$s. Let the amount of commencing $B$s be $j$, and the amount of terminating $R$s be $k$, I claim that the number of arcs which cross $(1,0)$ is constant, and it is equal to $\text{max}(j,k)$. I will show this with induction.

 Base case is when $n=1$. In this case, there are only two possible sequences - $RB$ and $BR$. In the first case, $\stackrel{\frown}{R_1B_1}$ does not cross $(1, 0)$, but both $j$ and $k$ are $0$, so $\text{max}(j,k)=0$. In the second example, $j=1$, $k=1$, so $\text{max}(j,k)=1$. $\stackrel{\frown}{R_1B_1}$ crosses $(1,0)$ since $B_1$ appears to the left of $R_1$, so there is one arc which intersects. Hence, the base case is proved.

 For the inductive step, suppose that for a positive number $n$, the number of arcs which cross $(1,0)$ is constant, and given by $\text{max}(j, k)$ for any configuration. Now, I will show it for $n+1$.

 Suppose I first choose $R_1$ such that $B_1$ is to the right of $R_1$ in the sequence. This implies that $\stackrel{\frown}{R_1B_1}$ does not cross $(1,0)$. But, neither $R_1$ nor $B_1$ is a commencing $B$ or terminating $R$. These numbers remain constant, and now after this move we have a sequence of length $2n$. Hence, by assumption, the total amount of arcs is $0+\text{max}(j,k)=\text{max}(j,k)$.

 
                \begin{itemize}
                    \item Here is a counter-case. $BRR_{1}B_{1}RR$ : j = 1, k = 2 => $BRRR$ : j = 1, k = 3. These numbers may not remain constant.
                \end{itemize}


                Thus, this solution probably doesn't work.

 Now suppose that $R_1$ appears to the right of $B_1$, but $B_1$ is not a commencing $B$. This implies that there are no commencing $B$s in the series, because there are no $B$s to the left of $B_1$, so $j=0$. Note that this arc does intersect $(1,0)$, and $R_1$ must be a terminating $R$. $R_1$ must be a terminating $R$ because there are no $B$s to the right of $R_1$, or else that $B$ would be $B_1$. The $2n$ length sequence that remains has $0$ commencing $B$s and $k-1$ terminating $R$s. Hence, by assumption, the total amount of arcs is $1+\text{max}(0,k-1)=1+k-1=k=\text{max}(j,k)$.

 Finally, suppose that $R_1$ appears to the right of $B_1$, and $B_1$ is a commencing $B$. We know that this arc will cross $(1,0)$. Analogous to the previous case, $R_1$ is a terminating $R$, so the $2n$ length sequence which remains has $j-1$ commencing $B$s and $k-1$ terminating $R$s. Hence, by assumption, the total amount of arcs is $1+\text{max}(j-1,k-1)=1+\text{max}(j,k)-1=\text{max}(j,k)$.

 There are no more possible cases, hence the induction is complete, and the number of arcs which intersect $(1,0)$ is indeed a constant which is given by $\text{max}(j,k)$.

 -william122

 



\textbf{Solution 2}

Lemma: If we switch the ordering of two consecutive $R_k$, $R_{k+1}$, the number of arcs crossing $(1,0)$ stays invariant.

 Proof: There are two situations. If the two arcs don't cross this is simple because the actual arcs stay the same, and only the number order of the arcs change. Otherwise, if $B_k$ is to the left of $R_{k+1}$, the two arcs will have some overlap. The order of the points counterclockwise must be RRBB or some rotation of it. Notice how in both the old and the new ordering, the arc between the second R and the first B is counted twice and the arcs between the same colors are only counted once. Thus no matter where $(1,0)$ is the number of arcs containing it will stay invariant.

 Lemma 2: We can swap the place of any two points $R_i$ and $R_j$ while keeping the number of crossings invariant.

 Proof: By Lemma 1 we can use consecutive arc moves to swap $R_i$ with $R_{i+1}$, then the new $R_{i+1}$ with $R_{i+2}$ and so on until it swaps with $R_j$. Now, we can do the reverse swaps to move the new $R_{j-1}$ to $R_{j-2}$, and so on into the original position of $R_i$. Since the other points were moved once in one direction and once in the opposite, all the other points are kept in place.

 Now, let $I$ be the sequence where $R_1$ through $R_n$ is numbered counterclockwise. We show that any ordering has the same number of crossings as $I$. We can swap the current place of $R_1$ with the desired position in $I$. Now we can ignore $R_1$ and swap the currrent place of $R_2$ with the place of $R_2$ in $I$, and so on until all the points are in the desired position. Thus, all orderings of $R_1$ to $R_n$ have the same number of crossings as $I$.

 -tigershark22

 


	\item (\textit{USAMO 2017}) Note: For any geometry problem whose statement begins with an asterisk ($*$), the first page of the solution must be a large, in-scale, clearly labeled diagram. Failure to meet this requirement will result in an automatic 1-point deduction. 


 Let $\mathbf{Z}$ denote the set of all integers. Find all real numbers $c > 0$ such that there exists a labeling of the lattice points $( x, y ) \in \mathbf{Z}^2$ with positive integers for which:
only finitely many distinct labels occur, and
for each label $i$, the distance between any two points labeled $i$ is at least $c^i$.


 



\textbf{Solution 1 (Evan Chen, Calvin Deng)}

See page 11 of this PDF:\href{https://web.evanchen.cc/exams/USAMO-2017-notes.pdf}{https://web.evanchen.cc/exams/USAMO-2017-notes.pdf} 

 


 



\textbf{Solution 2 (INCOMPLETE)}

For $c\le 1,$ we can label every lattice point $1.$ For $c\le \sqrt[4]{2},$ we can make a "checkerboard" labeling, i.e. label $(x, y)$ with $1$ if $x+y$ is even and $2$ if $x+y$ is odd. One can easily verify that these labelings satisfy the required conditions. Therefore, a labeling as desired exists for all $0 < c\le \sqrt[4]{2}.$

 An iterated version of the checkerboard labeling can actually work for all values $c < \sqrt{2}.$ For convenience, define the original lattice grid to be the set of all lattice points in the coordinate plane. Define a modified lattice grid of size $x$ to be a structure similar to the lattice points on the coordinate plane, but with the minimum separation between any two points equaling $x$ (as opposed to $1$).

 On the first step, assign a label $1$ to half of the points in a checkerboard arrangement. One can see that the points that have not yet been labeled form a modified lattice grid of size $\sqrt{2}$ (this lattice grid is also rotated by $45^{\circ}$ from the original lattice grid). At this point, for the second step, assign a label $2$ to half of the points, again in a checkerboard arrangement. At this point, the points that have not yet been labeled form a modified lattice grid of size $2$ (and again, it is rotated $45^{\circ}$ from the modified lattice grid after the first step). One then continues in this fashion. For the $N^{\text{th}}$ step, the points we are labeling are separated by at least $\sqrt{2}\times\left(\sqrt{2}\right)^{N-1} = \left(\sqrt{2}\right)^N > c^N,$ so we know that our labeling at each step is acceptable.

 After the $N^{\text{th}}$ step (where $N$ is a natural number), the points that have not yet been labeled form a modified lattice grid with size $\left(\sqrt{2}\right)^N.$ Since $c < \sqrt{2},$ we will eventually have $\left(\sqrt{2}\right)^N > c^{N+1}$ for some sufficiently large $N.$ At this point, we can label all remaining points in the original lattice grid $N+1,$ and this produces a labeling of all of the lattice points in the plane that satisfies all of the conditions. Therefore, a labeling as desired exists for all $c < \sqrt{2}.$

 We now prove that no labeling as desired exists for any $c\ge 2.$ To do this, we will prove that labeling a $2^k$-by-$2^k$ square grid of lattice points requires at least $k+3$ distinct labels for all natural numbers $k$; hence for a sufficiently large section of the lattice plane the number of distinct labels required grows arbitrarily large, so the entire lattice plane cannot be labeled with finitely many distinct labels. We will prove this using induction.

 For the base case, $k=1,$ we have four points in a square of side length $1.$ The maximum distance between any two of these points is $\sqrt{2} < c^1$ for all $c\ge 2,$ so all four points must have different labels. This completes the base case.

 Now, for the inductive step, suppose that labeling a $2^k$-by-$2^k$ square grid of lattice points requires at least $k+3$ distinct labels for some natural number $k.$ We will now prove that labeling a $2^{k+1}$-by-$2^{k+1}$ square grid of lattice points requires at least $k+4$ distinct labels.

 Take a $2^{k+1}$-by-$2^{k+1}$ square grid of lattice points. Divide this grid into four quadrants, $A, B, C,$ and $D.$ By the inductive hypothesis, $A$ requires at least $k+3$ distinct labels. At least one of these labels must be $k+3$ or greater; take one such label and call it $L.$

 The largest distance between any two points in the entire grid is $\sqrt{2}\left(2^{k+1} - 1\right) < c^{k+3}$ for all $c\ge 2.$ Therefore, the label $L$ cannot be used anywhere else in the grid. However, $B, C,$ and $D$ each require at least $k+3$ distinct labels as well by the inductive hypothesis. Thus, they must use at least one label that is not used in $A.$ It follows that the entire grid requires at least $k+4$ distinct labels. This completes the inductive step, and thus we conclude that no labeling as desired exists for any $c\ge 2.$

 I have heard from others that the actual boundary is $\sqrt{2}.$ This makes intuitive sense, since the iterated checkerboard labeling outlined above just breaks down at this value (you will be able to get closer and closer to labeling all of the lattice points, but you can never get there, since you will never have $\left(\sqrt{2}\right)^N > c^{N+1}$). The inductive argument above seems fairly loose, so I think that it can be sharpened to bring the upper bound down to $\sqrt{2},$ but I am not sure yet how exactly to do so. I think the way to do it is to somehow force $2$ new labels (instead of just $1$) each time you double the side length of the square grid.

 



\textbf{Solution 3}

We claim that for all $c<\sqrt{2}$, there exists such a configuration. 

 Lemma 1a: If there exists a configuration for $c=c_1$, then for any $c_2\le{c_1}$, there also exists a configuration for $c=c_2$.Note that any labeling that can be done for $c=c_1$ exists for $c=c_2$. Thus, we may copy the configuration from $c=c_1$ to $c=c_2$. 

 Lemma 1b: If there does not exist a configuration for $c=c_1$, then for any $c_2\ge{c_1}$, there does not exist a configuration for $c=c_2$.Note that any labeling which may be done for $c=c_2$ should also be available for $c=c_1$. However, we already stated that $c=c_1$ has no configuration, and so $c=c_2$ doesn't either. 

 We now prove that $c=\sqrt{2}$ does not work. WLOG, label the origin of our lattice plan with label $1$. WLOG, we can cover every point that we can with label $1$ because not labeling any such possible point does not affect the distance between points that aren't allowed to be labelled. 

 Name the state of the lattice grid before using label $1$ state $S$. After covering every point possible, we are left with only lattice points that have coordinates that sum to an odd integer. Notice, however, that if we discard the covered points and rotate all the new points by $\frac{\pi}{4}$ radians counterclockwise, we end up with the same grid as state $S$, only all the distances are multiplied by scalar $\sqrt{2}$. However, before placing label $2$s, note that we are in the same state $S$ as before, because our minimum distance is also multiplied by scalar $\sqrt{2}$. Thus, we will always be in an infinite loop, stuck in state $S$, and thus it is impossible to complete the labeling. 

 It suffices to show that any $c=2^{k}$ works given $k<0.5$. Note that the two smallest distances between any two points in the starting grid are $1$ and $\sqrt{2}$. When we are currently in state $S$ using label $I$, it is obvious that if $2ki-i\le{1}$, we can break out of state $S$ (Our effective minimum length would be at most 1, and we can cover every remaining point). Because $k<0.5$, there does exist an $I$ where this occurs. Thus, any $c<\sqrt{2}$ works, hence proved. $\blacksquare{}$ ~SigmaPiE

 


	\item (\textit{USAMO 2017}) Note: For any geometry problem whose statement begins with an asterisk ($*$), the first page of the solution must be a large, in-scale, clearly labeled diagram. Failure to meet this requirement will result in an automatic 1-point deduction. 


 Find the minimum possible value of 
 \[\frac{a}{b^3+4}+\frac{b}{c^3+4}+\frac{c}{d^3+4}+\frac{d}{a^3+4}\]given that $a$, $b$, $c$, $d$ are nonnegative real numbers such that $a+b+c+d=4$.


 



\textbf{Solution}

See here: \href{https://artofproblemsolving.com/community/c5t211539f5h1434574\_looks\_like\_mount\_inequality\_erupted\_}{https://artofproblemsolving.com/community/c5t211539f5h1434574\_looks\_like\_mount\_inequality\_erupted\_}

 or:

 


\end{enumerate}